\chapter*{Marking Scheme}
\addcontentsline{toc}{chapter}{Marking Scheme}

This document is structured to address all marking criteria for the Software Engineering G6046 group coursework. The table below outlines each requirement, the maximum marks available, and where evidence of meeting that requirement can be found within this report.

\vspace{1cm}

\begin{center}
\begin{tabular}{|p{3.2cm}|p{1.5cm}|p{5.8cm}|p{2.2cm}|}
\hline
\rowcolor{darkblue}
\textcolor{white}{\textbf{Element}} &
\textcolor{white}{\textbf{Marks}} &
\textcolor{white}{\textbf{Description}} &
\textcolor{white}{\textbf{Link}} \\
\hline

Planning Docs (out of 5) & 5 & Project plan, PERT/Gantt charts, risk analysis & \hyperref[chap:planning]{\color{blue}\underline{Chapter \ref*{chap:planning}}} \\
\hline

Planning --- Meeting Notes (out of 5) & 5 & Meeting records, action tracking, team decisions & \hyperref[sect:meeting]{\color{blue}\underline{Section \ref*{sect:meeting}}} \\
\hline

Process Documentation (out of 10) & 10 & Sprint documentation, user stories, Agile evidence & \hyperref[chap:process]{\color{blue}\underline{Chapter \ref*{chap:process}}} \\
\hline

Design Documentation --- High Level (out of 5) & 5 & System architecture, component roles, interactions & \hyperref[chap:design-high]{\color{blue}\underline{Chapter \ref*{chap:design-high}}} \\
\hline

Design Documentation --- Low Level (out of 5) & 5 & Class diagrams, OO principles, API design & \hyperref[chap:design-low]{\color{blue}\underline{Chapter \ref*{chap:design-low}}} \\
\hline

Software Documentation (out of 5) & 5 & Code comments, XML DocStrings, IntelliSense cross-refs & \hyperref[chap:software]{\color{blue}\underline{Chapter \ref*{chap:software}}} \\
\hline

Testing Documentation (out of 10) & 10 & Unit tests, system tests, requirements traceability, edge cases & \hyperref[chap:process]{\color{blue}\underline{Chapter \ref*{chap:process}}} \\
\hline

Code (out of 50) & 50 & Working game, AI player, GUI, customisation & Submitted codebase \\
\hline

Group Report (out of 5) & 5 & Achievements, issues, critical analysis, team organisation & \hyperref[chap:project-overview]{\color{blue}\underline{Chapter \ref*{chap:project-overview}}} \\
\hline

\rowcolor{lightblue}
\textbf{TOTAL} & \textbf{100} & \textbf{Team Mark Base} & --- \\
\hline

\end{tabular}
\end{center}

\vspace{1em}

\noindent\textbf{Known gaps against the highest marking bands.} The following limitations prevent a full claim on the top band for some elements and are acknowledged explicitly in the report:

\begin{itemize}
    \item \textbf{Testing (9--10 band):} Sprint~4 system tests (S4.1--S4.12) were not fully executed before submission --- the Actual Result and Pass/Fail columns remain incomplete (Section~\ref{sec:sprint4}). Edge-case tests were performed for board movement (Sprint~3) but are not formally catalogued as a dedicated edge-case suite. Unit tests are linked to classes but not formally cross-referenced back to individual design artefacts (class diagram entries or sequence diagrams).
    \item \textbf{Code (41--50 band):} The AI player uses BFS-directed navigation toward envelope-candidate rooms and knowledge-filtered suggestions, but full notebook-based card-elimination deduction was not implemented (D3, KI-3 in Section~\ref{sec:sprint4-known-issues}). The player-count selector in the menu is not fully wired to \texttt{SceneController} (KI-1), and player names are not user-configurable --- Watson Games specifically identified customisation as a selling point. The game does not feature animated character movement or formally verified cross-platform (Mac/PC) operation.
\end{itemize}

\newpage
